% RESULTADO_RECUPERADO: cuerpo íntegro de
% output/CH_REFLEXION_REVISION_QUIRURGICA_20260917/ES/sections/ch_reflexion_testigos.tex.
% Cambios editoriales en el cuerpo recuperado: namespace ix: y jerarquía del título.
% La subsección final formaliza el empalme editorial con el consumidor X.
% Residencia: inmediatamente después del lema ix:thm:ch-6b.
\subsection{Réflexion par clôture de témoins et construction de l'image}
\label{empalme:reflexion-testigos}
L'étape de Collection--Remplacement du lemme~\ref{ix:thm:ch-6b} admet la
démonstration explicite suivante. On maintient la hiérarchie déjà engendrée à partir du
code opératoire et du registre réalisé ; nous abrégeons
\(\mathfrak G=\mathfrak G_{\rm HMT}(h,m_\infty)\).
L'argumentation utilise Séparation et Remplacement dans la métathéorie ambiante,
sans présupposer ces schémas dans \(\mathfrak G\).

Soit \(\Phi\) une collection finie de formules. On traduit d'abord leurs quantificateurs universels au moyen de la négation et de la quantification existentielle, puis on ferme la collection résultante par sous-formules. Pour tout niveau initial, il existe un ordinal limite ultérieur \(\theta\) tel que
\[
 \mathfrak G_\theta\models\varphi(\bar a)
 \quad\Longleftrightarrow\quad
 \mathfrak G\models\varphi(\bar a)
 \qquad
 (\varphi\in\Phi,\ \bar a\in\mathfrak G_\theta^{<\omega}).
\]
La notation de satisfaction de la classe est interprétée, pour chaque formule fixée, en relativisant ses quantificateurs à la classe définissable \(\mathfrak G\) ; aucun prédicat uniforme de vérité n'est introduit.

\begin{proof}
Fixons \(\alpha\). Pour chaque sous-formule existentielle
\(\exists y\,\psi(y,\bar x)\) de \(\Phi\), les tuples
\(\bar a\in\mathfrak G_\alpha^{<\omega}\) de l’arité correspondante
qui possèdent un témoin dans \(\mathfrak G\) forment un ensemble par Séparation
métathéorique. À chacun est attribué le plus petit ordinal \(\beta\) pour lequel
un \(y\in\mathfrak G_\beta\) satisfait
\(\mathfrak G\models\psi(y,\bar a)\). L’ordinal existe car
\(\mathfrak G\) est l’union de ses niveaux. Le Remplacement dans la métathéorie
borne ces premières étapes des témoins. La finitude de \(\Phi\) permet
de choisir de manière définissable une seule borne \(f(\alpha)>\alpha\), valable pour
toutes ses sous-formules existentielles. Par conséquent, chaque tuple considéré possède
\emph{au moins un témoin} dans \(\mathfrak G_{f(\alpha)}\).

Nous prenons \(\alpha_0\) suffisamment grand pour contenir les paramètres
initiaux et posons
\[
 \alpha_{n+1}=f(\alpha_n),\qquad
 \theta=\sup_{n<\omega}\alpha_n.
\]
Par continuité de la hiérarchie, tout tuple fini de
\(\mathfrak G_\theta\) appartient à un certain \(\mathfrak G_{\alpha_n}\).
Si une formule existentielle de \(\Phi\) avec ce tuple est vraie dans
\(\mathfrak G\), elle possède un témoin dans
\(\mathfrak G_{\alpha_{n+1}}\subseteq\mathfrak G_\theta\).
La récurrence sur les formules prouve maintenant l’équivalence : les formules atomiques
sont conservées dans la structure induite ; les connecteurs préservent
l’équivalence ; l’étape existentielle utilise ce témoin dans un sens et le
même témoin, vu dans \(\mathfrak G\), dans l’autre. La traduction préalable
des universelles inclut aussi les formules existentielles qui expriment
leurs contre-exemples possibles. La réflexion finie requise est prouvée.
\end{proof}

Soit maintenant \(F(x,y,\bar p)\) fonctionnelle sur \(a\) dans \(\mathfrak G\). Nous appliquons la construction aux formules qui expriment \(F\), l'existence de ses témoins et son image, avec \(a,\bar p\in\mathfrak G_\theta\). La transitivité donne \(a\subseteq\mathfrak G_\theta\). Chaque \(x\in a\) possède sa valeur à l'intérieur de ce niveau, et la réflexion identifie l'image exacte à
\[
 b=\{y\in\mathfrak G_\theta:
       \mathfrak G_\theta\models
       \exists x\in a\ F(x,y,\bar p)\}.
\]
Cet ensemble est définissable sur \(\mathfrak G_\theta\), avec des paramètres du même niveau. Par l'opération successeur de la hiérarchie, \(b\in\operatorname{Def}(\mathfrak G_\theta)
\subseteq\mathfrak G_{\theta+1}\). On obtient ainsi le remplacement dans \(\mathfrak G\). Si l'on conserve l'existence de témoins et que l'on omet l'unicité, la même construction fournit un ensemble de témoins pour la collection.

\paragraph{Arité des noms et sélection interne.}
Dans le bon ordre du lemme~\ref{ix:thm:ch-6b}, l’énumération de la base
attribue à chaque élément son plus petit code naturel. À chaque successeur, on ordonne
d’abord les codes de formule ; une fois l’un fixé, l’arité de son
tuple de paramètres est fixée. L’ordre lexicographique s’applique alors à un
\emph{produit fini d’arité fixe} des bons ordres précédents.
Le plus petit nom représente chaque nouvelle valeur et les niveaux limites conservent
l’ordre de première apparition. L’ordre des tuples est ainsi spécifié,
sans mélanger des arités différentes dans un unique ordre
lexicographique. La sélection du plus petit élément de chaque ensemble non vide,
avec le Remplacement déjà obtenu, forme intérieurement les fonctions de
choix utilisées ensuite. Cette construction de noms précède les
injections cardinales et ne reçoit aucune bijection du continu comme
donnée.

\subsection{Compatibilité du classificateur cardinal avec la restitution des histoires}
\label{empalme:clasificador-restitucion}

La proposition~\ref{x_reciprocidad:lib07:prop-familias} construit, à partir de
lecteurs finis inversibles compatibles avec la troncature, une bijection
\(q:H\to Y\) entre la limite des histoires enrichies et la limite de leurs
registres récupérables. Son domaine conserve le système de niveaux, les
opérations et la mémoire. L'argument cardinal utilise en outre la lecture
archimédienne \(p:H\twoheadrightarrow R\), où
\(R=\mathbb R_{\rm HMT}^{\mathfrak G}\). Une histoire et sa valeur
publiée sont des objets distincts : les fibres de \(p\) rassemblent les
réalisations enrichies d'une même valeur.

\begin{proposition}[Transport interne des fibres et conservation du code ordinal]
\label{empalme:prop-fibras-codigo}
Fixons la clôture \(\mathfrak G\) de la section~\ref{ix:sec:cierre-cardinal-nativo}.
Supposons que les ensembles \(H,Y,R\), la bijection \(q\), la
surjection \(p\) et leurs systèmes définissants appartiennent à
\(\mathfrak G\). Définissons
\[
 \widetilde p=p\circ q^{-1}:Y\longrightarrow R,
 \qquad
 \mathcal S(A)=q[p^{-1}(A)]
 \quad(A\in\mathcal P^{\mathfrak G}(R)).
\]
Alors \(\mathcal S\) est un isomorphisme entre
\(\mathcal P^{\mathfrak G}(R)\) et l'algèbre des sous-ensembles internes
de \(Y\) saturés pour \(\widetilde p\). On a
\[
 \mathcal S(A)=\widetilde p^{-1}(A),\qquad
 \widetilde p[\mathcal S(A)]=A,
 \qquad
 (j_{\rm HMT}\circ\widetilde p)\circ q=j_{\rm HMT}\circ p.
\]
La restriction \(q:p^{-1}(A)\to\mathcal S(A)\) est bijective. Le
classificateur \(j_{\rm HMT}\) de~\eqref{ix:eq:ch-indices-real} conserve
son code lors du changement de registre récupérable et maintient la dichotomie
\eqref{ix:eq:ch-dicotomia} pour les valeurs de \(A\).
\end{proposition}

\begin{proof}
Séparation forme le graphe inverse de \(q\) dans \(Y\times H\).
Remplacement, déjà obtenu par la clôture de témoins, forme les compositions
et les images utilisées. Pour chaque \(y\in Y\), l'unicité de
\(q^{-1}(y)\) donne
\[
 y\in q[p^{-1}(A)]
 \ \Longleftrightarrow\ p(q^{-1}(y))\in A
 \ \Longleftrightarrow\ y\in\widetilde p^{-1}(A).
\]
La surjectivité de \(p\) et la bijection \(q\) prouvent
\(\widetilde p[\widetilde p^{-1}(A)]=A\). Si \(B\subseteq Y\)
est saturé, par définition
\(B=\widetilde p^{-1}(\widetilde p[B])\), si bien que \(B\)
appartient à l'image de \(\mathcal S\). Les images réciproques préservent
les unions, les intersections et les compléments relatifs ; la surjectivité de
\(\widetilde p\) rend cette opération injective. Cela prouve
l'isomorphisme des algèbres de sous-ensembles. La restriction de \(q\) est
bijective parce que \(q\) l'est sur son domaine complet.

Enfin, \(\widetilde p\circ q=p\), et la composition avec
\(j_{\rm HMT}\) donne le carré déclaré. Les fonctions \(e_\beta\),
la condensation et les codes ordinaux sont construits dans la même
clôture par les lemmes~\ref{ix:thm:ch-7} et~\ref{ix:thm:ch-8} ;
l'égalité~\eqref{ix:eq:ch-igualdad-cardinal} et la dichotomie
\eqref{ix:eq:ch-dicotomia} s'appliquent à \(A\) dans ce domaine. On
transporte ainsi la classification des valeurs en conservant toutes leurs
fibres enrichies.
\end{proof}

Le résultat distingue les deux bijections qui interviennent : la restitution
agit entre histoires et registres ; la classification cardinale agit sur
les valeurs publiées. La saturation relie les deux niveaux sans
sélectionner une seule histoire par valeur. La réflexion des témoins garantit
la formation interne des applications de l'argument et la conservation des
quantificateurs de la clôture déclarée.
